\documentclass[10pt,a4paper]{amsart}
\usepackage[margin=19mm]{geometry}
\usepackage{amsmath,amssymb,mathtools}
\usepackage[scr=rsfs,cal=euler]{mathalfa}
\usepackage{xfrac}
\usepackage[unicode,hidelinks,bookmarks=false]{hyperref}
\newcommand{\ie}{i.e.\ }
\newcommand{\cf}{cf.\ }
\newcommand{\resp}{respectively\ }
\newcommand{\Subsection}{\subsection}
\providecommand{\nobreakdash}{\nobreak\hbox{-}}
\setlength{\emergencystretch}{2em}
\multlinegap0pt
\usepackage{bm}
\usepackage{bbm}
\usepackage{dsfont}
\usepackage{xspace}
\usepackage{cancel}
\usepackage{mathtools}
\usepackage{cleveref}
\usepackage{amsthm}
\makeatletter
\@ifundefined{theorem}{\newtheorem{theorem}{Theorem}}{}
\@ifundefined{lemma}{\newtheorem{lemma}[theorem]{Lemma}}{}
\@ifundefined{proposition}{\newtheorem{proposition}[theorem]{Proposition}}{}
\makeatother
\newcommand*{\cA}{\mathcal{A}}
\newcommand*{\cB}{\mathcal{B}}
\newcommand*{\cC}{\mathcal{C}}
\newcommand*{\from}{\colon} 
\newcommand*{\gen}[1]{\left\langle #1\right\rangle}
\newcommand*{\lang}{\mathcal{L}}
\newcommand*{\ret}{\mathcal{R}}
\title[Stability properties for subgroups generated by return words]{Stability properties for subgroups generated by return words}
\author{France Gheeraert, Herman Goulet-Ouellet, Julien Leroy, Pierre Stas}
\date{Source: arXiv:2410.12534v1 (CC BY 4.0)}
\begin{document}
\maketitle

\noindent\emph{Source and licence.} Stability properties for subgroups generated by return words. Version arXiv:2410.12534v1 (CC BY 4.0), distributed under the Creative Commons Attribution 4.0 International licence, as recorded in the arXiv licence metadata for this exact version and shown on the abstract page https://arxiv.org/abs/2410.12534v1. Full rights notice: licenses/arxiv-2410.12534-CC-BY-4.0.txt.\par\smallskip

\noindent\textit{Editorial scope.} Counted unit: the Proposition whose source label is \texttt{P:eventual stability preserved by substitution} in the article (source0.tex lines 764--766, source label \texttt{P:eventual stability preserved by substitution}), delivered together with the article's own complete original proof of it (source0.tex lines 768--806, one \texttt{proof} environment of 39 lines). No mathematics is written, repaired or re-derived: statement and proof are the article's own text line for line. Required context retained uncounted: L:stable return groups also works for sets (lines 752--754). Every definition, display and label of those lines is retained unchanged. Omitted from this package and not redistributed: 1--751, 755--763, 767--767, 807--1042. Nothing inside a retained excerpt is omitted or altered: every retained body is delivered exactly as the article's lines are written, so no omission receipt of any kind accompanies this package. Citations: the retained excerpts carry no citation command at all, so no bibliography entry of the article is transcribed and no reference is redistributed. Reference adaptation: none was needed and none was made; every reference inside the delivered unit resolves inside it, so no unresolved reference or citation is produced. Printed slips kept literal: no printed word, symbol, hypothesis or number of the counted statement or of its proof was altered. Layout and class: the article's own class is \texttt{amsart} with packages including inputenc, fontenc, amsmath, amssymb, amsthm, xcolor, graphicx, doi, hyperref, csquotes, enumitem, booktabs, microtype, tikz; the wrapper is the lane's pinned \texttt{amsart} editorial wrapper at 10pt on A4, which restates the theorem-like environments the retained text needs and adds the article's own macro definitions that the retained text needs (\texttt{\textbackslash cA}, \texttt{\textbackslash cB}, \texttt{\textbackslash cC}, \texttt{\textbackslash from}, \texttt{\textbackslash gen}, \texttt{\textbackslash lang}, \texttt{\textbackslash ret}), with extra wrapper packages (bm, bbm, dsfont, xspace, cancel, mathtools, cleveref). The delivered numbering is the unit's own. Assets: no retained span contains a figure environment, a graphics-inclusion command, a PDF-page inclusion, a table or a TikZ picture; no image file is copied into this package, the package declares no asset dependency and no image file is redistributed with it. Licence: version arXiv:2410.12534v1 of this article is distributed under the Creative Commons Attribution 4.0 International licence, as recorded in the arXiv licence metadata for this exact version and shown on the abstract page https://arxiv.org/abs/2410.12534v1. A full rights notice is delivered at licenses/arxiv-2410.12534-CC-BY-4.0.txt. Characters: every retained excerpt is pure ASCII, so the delivered engine, XeLaTeX with its default Latin Modern text font, reports no missing character.
\begin{lemma}\label{L:stable return groups also works for sets}
    Let $X$ be a minimal shift space over $\cA$ and let $\varphi\from F_\cA \to G$ be a group morphism. Then $X$ has eventually $\varphi$-stable return groups of threshold at most $M$ if and only if, for all factor codes $T \subseteq \lang(X)$, we have $\varphi\gen{\ret_X(T)} = \varphi\gen{\ret_X(S)}$, where $S = \{u_{[0,\min\{|u|,M\})} \mid u \in T\}$.
\end{lemma}

\begin{proposition}\label{P:eventual stability preserved by substitution}
    Let $X$ be a minimal shift space over $\cA$ and let $\sigma\from \cA^* \to \cB^*$ be a substitution. Let $\varphi\from F_\cA \to G$ and $\psi\from F_\cB \to H$ be group morphisms satisfying $\ker(\varphi) \leq \ker(\psi\sigma)$. If $X$ has eventually $\varphi$-stable return groups of threshold $M$, then $\sigma[X]$ has eventually $\psi$-stable return groups of threshold at most $\max\{1,M|\sigma|\}$.
\end{proposition}

\begin{proof}
    First, recall that $|\sigma| = \max\{|\sigma(a)| \mid a \in \cA\}$ and observe that the hypothesis on the kernels is equivalent to $\bar\sigma = \psi\sigma\varphi^{-1}\from \varphi(F_\cA) \to H$ being a well-defined group morphism. This morphism satisfies $\bar\sigma\varphi = \psi\sigma$.
    
    To study $\sigma[X]$, we first factorize the substitution $\sigma$ into $\tau\theta$ as follows.
    Let $\cA = \{a_1, \dots, a_d\}$, $n_i = |\sigma(a_i)|$, and $\cC = \{a_{i,j} \mid 1 \leq i \leq d, 0 \leq j \leq |\sigma(a_i)|-1\}$ where $a_{i,j} \ne a_{i',j'}$ if $i \ne i'$ or $j \ne j'$. We define $\theta\from \cA^*\to \cC^*$ and $\tau\from \cC^*\to \cB^*$ by
    \[
        \theta(a_i) = a_{i,0} \cdots a_{i,n_i-1}\quad\text{and}\quad \tau(a_{i,j}) = \sigma(a_i)_j.
    \]
    They clearly satisfy $\sigma = \tau\theta$.
    
    We consider the intermediary shift space $\theta[X]$ and show that it is eventually $\psi\tau$-stable of threshold at most $\max\{1,M|\sigma|\}$. To do so, we first study the return group of any non-empty $v \in \lang(\theta[X])$. Let $a_{i,j}$ and $a_{k,l}$ respectively be the first and last letters of $v$. Define $p_v = a_{i,0}\cdots a_{i,j-1}$ and $s_v = a_{k,l+1}\cdots a_{k,n_k-1}$. Then $p_vvs_v = \theta(t_v)$ for a unique $t_v\in\lang(X)$. Since $v$ only appears as a factor of $\theta(t_v)$ and $\theta(t_v)$ only appears as the image of $t_v$, we have
    \begin{equation}\label{eq:ret-theta}
        \ret_{\theta[X]}(v) = p_v^{-1}\ret_{\theta[X]}(\theta(t_v))p_v = p_v^{-1}\theta(\ret_X(t_v))p_v.
    \end{equation}
    Now, take $w \in \lang(\theta[X])$ of length $\max\{1,M|\sigma|\}$ and $u \in w\cC^* \cap \lang(\theta[X])$. We prove that $\psi\tau\gen{\ret_{\theta[X]}(w)} = \psi\tau\gen{\ret_{\theta[X]}(u)}$. By \cref{eq:ret-theta}, we have
    \begin{align*}
        \psi\tau\gen{\ret_{\theta[X]}(w)}
        &= \psi\tau(p_w)^{-1}\psi\tau\theta\gen{\ret_X(t_w)}\psi\tau(p_w) \\
        &= \psi\tau(p_w)^{-1}\psi\sigma\gen{\ret_X(t_w)}\psi\tau(p_w) \\
        &= \psi\tau(p_w)^{-1}\bar\sigma\varphi\gen{\ret_X(t_w)}\psi\tau(p_w).
    \end{align*}
    Similarly, $\psi\tau\gen{\ret_{\theta[X]}(u)} = \psi\tau(p_u)^{-1}\bar\sigma\varphi\gen{\ret_X(t_u)}\psi\tau(p_u)$.
    As $p_u = p_w$ by definition of $u$ and $w$, it then suffices to show that $\varphi\gen{\ret_X(t_w)} = \varphi\gen{\ret_X(t_u)}$. By definition of $u$ and $w$ once again, $t_w$ is a prefix of $t_u$ and
    \[
        |t_w| \geq \frac{|\theta(t_w)|}{|\sigma|} = \frac{|p_wwq_w|}{|\sigma|} \geq \frac{|w|}{|\sigma|}
        \geq M.
    \]
    As $X$ is eventually $\varphi$-stable of threshold $M$, this directly implies that $\varphi\gen{\ret_X(t_w)} = \varphi\gen{\ret_X(t_u)}$, and therefore that $\psi\tau\gen{\ret_{\theta[X]}(w)} = \psi\tau\gen{\ret_{\theta[X]}(u)}$. This ends the proof that $\theta[X]$ is eventually $\psi\tau$-stable of threshold $K \leq \max\{1,M|\sigma|\}$.
    
    Let us now show that $\sigma[X]$ is eventually $\psi$-stable of threshold at most $K$. Let $w \in \lang_K(\sigma[X])$ and let $u \in w\cB^* \cap \lang(\sigma[X])$. Since $\tau$ is letter-to-letter, one easily sees that $\tau^{-1}(u)\cap\lang(\theta[X])$ is a factor code and 
    \[
        \ret_{\sigma[X]}(u) = \tau(\ret_{\theta[X]}(\tau^{-1}(u)\cap\lang(\theta[X]))).
    \]
    Set $T = \tau^{-1}(u)\cap\lang(\theta[X])$ and $S = \{v_{[0,K)} \mid v \in T\}$. Note that $S = \tau^{-1}(w)\cap\lang(\theta[X])$ thus we similarly have $\ret_{\sigma[X]}(w) = \tau(\ret_{\theta[X]}(S))$. As $\theta[X]$ is eventually $\psi\tau$-stable of threshold $K$, we obtain by \cref{L:stable return groups also works for sets}
    \[
        \psi\ret_{\sigma[X]}(u) = \psi\tau(\ret_{\theta[X]}(T)) = \psi\tau(\ret_{\theta[X]}(S)) = \psi\ret_{\sigma[X]}(w).
    \]
    This shows that $\sigma[X]$ is eventually $\psi$-stable of threshold at most $K \leq \max\{1,M|\sigma|\}$.
\end{proof}

\end{document}
